\part{Valuations, models, and analytic spaces}
\label{part:foundations}

This part builds the common language of the book.  Its purpose is not only to
introduce Berkovich spaces, but to make four descriptions of the same basic
geometric data usable side by side: seminorms, valuations, models, and metric
weights.  The four entrances are displayed in \Cref{chap:four-entrances}.
Valuations and their centers are developed in \Cref{chap:valuations}; schemes
over valuation rings and models occupy \Cref{chap:models}; and the passage to
Berkovich spaces, skeleta, and metrics is made in
\Cref{chap:berkovich,chap:skeleta,chap:metrics}.  The topological and
measure-theoretic toolkit in \Cref{chap:topological-tools} prevents later
compactness and monotone-convergence arguments from silently replacing nets
by sequences.  Finally, \Cref{chap:model-calculations} turns curves and toric
varieties into a test suite for the signs, slopes, and masses used later.

\paragraph{Prerequisites and support.}
\label{front:part1-prerequisites-and-support}
A first course in schemes or compact complex geometry is enough to begin.
The chapters recall the valuation-theoretic and analytic definitions that
they use, while general facts about proper morphisms, normalization, and
intersection theory are cited at their point of application.  Readers who
want those algebraic inputs collected in one place may consult
\Cref{app:valuation-algebra,app:birational-geometry,app:stacks}; the global
sign and scaling conventions are gathered in \Cref{app:normalizations}.
No prior Berkovich geometry is assumed.

\paragraph{Conceptual clusters.}
\label{front:part1-conceptual-clusters}
The first cluster,
\Cref{chap:four-entrances,chap:valuations,chap:berkovich}, explains what an
analytic point records and how centers, residue fields, and functoriality are
read from it.  The model cluster,
\Cref{chap:models,chap:skeleta,chap:metrics}, passes from a proper model to its
dual complex, from a vertical Cartier divisor to a piecewise affine metric
weight, and from relative intersection theory to a model measure.  The final
cluster, \Cref{chap:topological-tools,chap:model-calculations}, supplies the
compactness, envelope, weak-measure, and convex tools used throughout the
book and then checks them on metric graphs, projective lines, and toric
skeleta.

\paragraph{Logical spine.}
\label{front:part1-logical-spine}
The recurring chain is
\[
\begin{aligned}
 \text{valuation}
 &\longrightarrow \text{analytic point}
 \longrightarrow \text{model and skeleton},\\
 &\longrightarrow \text{model metric}
 \longrightarrow \text{atomic intersection measure}.
\end{aligned}
\]
Each arrow comes with a compatibility test.  Centers are checked against the
valuative criterion of properness; skeleta against monomial coordinates;
model functions against vertical Cartier divisors; and measures against the
intersection number of the underlying line bundle.  These tests form the
common model locus on which the BFJ, CLD, and BJ constructions can later be
compared.

\paragraph{Reading routes.}
\label{front:part1-reading-routes}
The \emph{short route} is
\Cref{chap:four-entrances,chap:valuations,chap:berkovich,chap:metrics},
together with the curve and toric tests in \Cref{chap:model-calculations}; it
gives a working vocabulary for sampling the rest of the book.  The
\emph{standard route} reads the part in order and is the appropriate
preparation for \Cref{part:bfj,part:bj}.  The \emph{reference route} adds
\Cref{chap:topological-tools} and the supporting appendices
\Cref{app:valuation-algebra,app:birational-geometry,app:nets,app:intersection-bdivisors};
it is designed for readers who want to check every later compactness or
intersection-theoretic dependency.  A reader heading directly to the local
CLD calculus may postpone most of \Cref{chap:models,chap:skeleta}, but should
not postpone the analytic-space and metric conventions in
\Cref{chap:berkovich,chap:metrics}.

\paragraph{Main outcomes and interfaces.}
\label{front:part1-main-outcomes-and-interfaces}
At the end of the part, the reader should be able to move without ambiguity
between a multiplicative seminorm and its additive semivaluation, normalize a
divisorial point, construct the dual complex of an snc model, read a vertical
divisor as a piecewise affine function, and recover the total mass of a model
measure from intersection theory.  These operations feed the differential
theory of \Cref{part:cld} and the global BFJ theory of \Cref{part:bfj}.
After the ground norm becomes trivial, the same valuation language acquires
a scaling action and leads instead to \Cref{part:bj}.  In
\Cref{part:complex}, valuations and dual complexes reappear as the target of
complex singularity data; no identification across these settings is made
without the comparison results of \Cref{part:comparisons}.

\paragraph{Normalization warning.}
\label{front:part1-normalization-warning}
Throughout the book \(v=-\log|\cdot|\), effective vertical divisors produce
nonnegative model functions. Model and full-mass Monge--Amp\`ere measures
have their stated intersection-theoretic total mass.  Divisorial points over a nontrivially
valued field are normalized relative to the ground valuation, whereas points
on a trivially valued cone are compared only after a slice or scaling has
been specified.  The special-fibre, metric-weight, and factor-of-two tests in
\Cref{app:normalizations} should be consulted whenever two later formalisms
appear to differ only by a sign or constant.
